\documentclass{article}
\usepackage[utf8]{inputenc}
\usepackage{amsmath}
\usepackage{geometry}
\geometry{margin=2cm}
\begin{document}

\section*{Análise da Questão 152 — ENEM 2024 — Matemática}

\subsection*{Enunciado}
O gráfico apresenta o valor total de exportações e o valor total de importações, ao longo de um período, em bilhões de dólares. O saldo da balança comercial brasileira é dado pela diferença entre exportações e importações no mesmo período.

Os saldos destacados no gráfico para os meses são:
\begin{itemize}
  \item S_1: saldo em junho de 2009;
  \item S_2: saldo em janeiro de 2010;
  \item S_3: saldo em junho de 2010.
\end{itemize}

\subsection*{Solução Técnica}
\textbf{Fórmula utilizada:}
\[
\text{Saldo} = \text{Exportações} - \text{Importações}.
\]

Identificação dos valores aproximados no gráfico (em bilhões de dólares):
\[
\begin{cases}
S_1: (14,4, 11,0), \\
S_2: (11,5, 11,7), \\
S_3: (17,1, 14,8).
\end{cases}
\]

Cálculo dos saldos:
\[
\begin{aligned}
S_1 &= 14,4 - 11,0 = 3,4, \\
S_2 &= 11,5 - 11,7 = -0,2, \\
S_3 &= 17,1 - 14,8 = 2,3.
\end{aligned}
\]

Ordenando do maior para o menor:
\[
S_1 > S_3 > S_2.
\]

Portanto, a alternativa correta é \textbf{A}.

\subsection*{Explicação Didática}
Para determinar o saldo maior, calcula-se a diferença entre exportações e importações de cada mês e compara-se esse resultado.

\subsubsection*{Passos para resolução}
\begin{enumerate}
  \item Leia o gráfico e extraia valores aproximados de exportações e importações para os meses indicados.
  \item Calcule o saldo de cada mês usando a fórmula.
  \item Compare os saldos encontrados e ordene-os.
\end{enumerate}

\subsubsection*{Erros comuns}
\begin{itemize}
  \item Ignorar as importações no cálculo do saldo.
  \item Confundir os meses ou os valores no gráfico.
  \item Não realizar a subtração para descobrir o saldo real.
\end{itemize}

\subsubsection*{Dicas para o ensino}
\begin{itemize}
  \item Explique claramente o conceito do saldo da balança comercial.
  \item Mostre como interpretar gráficos para extrair dados.
  \item Enfatize a operação de subtração para obter o saldo.
  \item Promova vários exercícios de interpretação gráfica para fixação.
\end{itemize}

\subsection*{Distratores da Questão}
\begin{description}
  \item[Alternativa B:] Ignorar as importações, classificando de forma incorreta os saldos.
  \item[Alternativa C:] Considerar visualmente valores próximos como maiores sem calcular o saldo.
  \item[Alternativa D:] Valorizar somente as exportações altas sem considerar importações elevadas.
  \item[Alternativa E:] Inverter a ordem dos saldos baseando-se em leitura incorreta.
\end{description}

\subsection*{Perfil do Item}
\begin{itemize}
  \item Habilidade principal: interpretação e cálculo de dados em gráficos.
  \item Habilidades secundárias: análise quantitativa, diferença e ordenação de valores.
  \item Demanda cognitiva: médio.
  \item Dificuldade: média.
  \item Necessita interpretação visual, modelagem de subtração e comparação de resultados.
\end{itemize}

\subsection*{Revisão da Consistência}
A análise técnica e explicação didática estão coerentes com o enunciado e os dados do gráfico, com justificativas claras e resposta correta bem fundamentada. Os distractores são verossímeis e adequados para o nível da questão.

\end{document}